%titolo{Buon volo}
%autore{}
%album{}
%tonalita{DO}
%famiglia{Scout}
%gruppo{}
%momenti{}
%identificatore{buon_volo}
%data_revisione{2026_07_1}
%trascrittore{Modanese Tobia}

\beginsong{Buon volo}

\beginverse
\[C]Vola, \[F]coccinella, \[C]va', 
\[C]il viaggio è \[F]cominciato \[C]già 
\[F]per il bosco e per il \[C]prato in fior 
\[F]fino al monte dov'è l'\[C]aquila.
\[C]Con le for\[F]miche impari \[C]che 
\[C]tutte in\[F]sieme si sta \[C]ben: 
\[F]se ciascuno il meglio \[C]dà di sé
\[F]sempre il meglio noi pos\[G]siamo \[C]aver. 
\endverse

\beginverse
\[A-]Al girino tu fa\[E-]rai capir \[A-]che le zampe devo\[E-]no servir 
\[E]e così la rana \[A-]alfine sa \[F]saltellare qui e \[G]là. 
\endverse

\beginchorus
\[C]Coccinella buon \[G]volo, 
\[C]Coccinella buon \[G]volo, buon volo! 
\[C]Prendi un puntino in \[G]più, 
\[C]dai, buttati anche \[G]tu! 
\endchorus
\chordsoff
\beginverse
Vola, Coccinella, va', 
la civetta può parlar:
tutto è suo e se lo può tener,
noi vogliamo dare e non aver.
Gli scoiattoli vedrai 
giocare insieme e lavorar;
bello è stare nella quercia, ma
il tuo viaggio non finisce qua.
\endverse

\beginverse
Se il tuo volo non continuerai,
tu felice non lo sarai mai
e il tuo impegno ti richiama là:
la montagna vedi già. 
\endverse

\beginverse
Vola, Coccinella, va', 
il viaggio presto finirà
per il bosco e per il prato in fior
fino al monte dov'è l'aquila.
Il tuo aiuto non negar 
al serpente mai perchè
la risposta è una fatica, ma
dimmi cosa posso far per te.
\endverse

\beginverse
Occhi attenti sempre a scoprir
nuove cose a non finir
e per questo sai vedere chi 
al tuo aiuto dice sì.
\endverse

\endsong
